\section{Names and Folding}
\label{sec:names}

\subsection{The name equivalence}

The rules are matched against names that the operating system will treat as
identical.  On a case-insensitive APFS volume two names denote the same file
exactly when they are equal after Unicode normalization and \emph{full} case
folding.  Full folding differs from simple folding: it expands some characters
into several.

\begin{definition}[Fold]
\label{def:fold}
$\fold{s}=\Nfc\bigl(\mathsf{cf}(\Nfc(s))\bigr)$, where $\mathsf{cf}$ is Unicode
full case folding.  For names, $x\approx y \iff \fold{x}=\fold{y}$.  The
extension to paths applies $\Fold$ to the whole string, which is equivalent to
folding component by component because folding never introduces or removes a
separator.
\end{definition}

\begin{law}[Idempotence]
\label{law:fold-idempotent}
$\fold{\fold{s}}=\fold{s}$.
\end{law}

\begin{observation}[Measured aliases]
\label{obs:aliases}
The set of single characters that APFS folds onto ASCII names was measured by
creating every ASCII name of length one and two, and then probing every
character in the Unicode range \code{U+0080}--\code{U+2FFFF}.  On the tested
system the aliases are exactly the following; each row names the ASCII
spelling that the character opens.
\begin{center}
\begin{tabular}{@{}ll@{}}
\toprule
character & opens \\
\midrule
\code{\ss} (sharp s), capital sharp s (\code{U+1E9E}) & \code{ss} \\
long s (\code{U+017F}) & \code{s} \\
Kelvin sign (\code{U+212A}) & \code{k} \\
ligatures \code{U+FB00}--\code{U+FB06} & \code{ff}, \code{fi}, \code{fl}, \code{ffi}, \code{ffl}, \code{st} \\
\bottomrule
\end{tabular}
\end{center}
This is full case folding, not simple folding: under simple folding
\code{\ss} would remain a distinct character.
\end{observation}

\begin{law}[Fold agrees on the measured aliases]
\label{law:fold-aliases}
For every pair in Observation~\ref{obs:aliases},
$\fold{\text{character}}=\fold{\text{ASCII spelling}}$.  The server's $\Fold$
(Go, \code{x/text/cases}) and the client's $\Fold$ (JavaScript,
lower-, upper-, then lower-casing) both satisfy the law.
\Tests{rules/fold_test.go::TestFoldMatchesTheAliasesAPFSTreatsAsEqual; web/src/lib/laws.test.ts (filter agreement)}
\end{law}

\begin{boundary}[Folding is a property of the volume]
\label{bnd:volume}
The alias set is a fact about the file system, not about \texttt{fsb}.  A
case-sensitive volume distinguishes more names and a future OS release could
change the table.  The specification therefore treats $\approx$ as the
\emph{coarsest} equivalence the program must respect: comparing after folding
can only over-match, and over-matching a deny rule fails closed.
\end{boundary}
